\documentclass{article}
\usepackage[utf8]{inputenc}
\usepackage{amsmath,amssymb}
\usepackage[most]{tcolorbox}
\usepackage{geometry}
\geometry{margin=2.5cm}

\newcommand{\step}[1]{\par\noindent\hangindent=3.5em\hangafter=1%
  \makebox[3.5em][l]{\textbf{#1.}}}
\newcommand{\steptag}[1]{\hfill{\small\textsf{[#1]}}}

\newtcolorbox{proofbox}[1]{
  colback=gray!5, colframe=gray!60,
  fonttitle=\bfseries, title={#1},
  breakable, enhanced,
  left=4pt, right=4pt, top=4pt, bottom=4pt,
  before skip=10pt, after skip=10pt
}

\begin{document}

\begin{proofbox}{After first fixing one global witness package W\_RF suppli...}
\step{1} After first fixing one global witness package W\_RF supplied by lem-routef-raw-factor-setting-formation, for every input (H,Phi,eta) to which that formation result applies, fix one def-routef-raw-factor-setting datum S over that same W\_RF supplied by the same result, writing the fields of (W\_RF,S) as the unqualified symbols below: for every Delta' supplied for (W\_RF,S) by lem-routef-delta-prime-closeness and every Delta supplied from that same Delta' by lem-routef-delta-normalization-closeness, with C\_N, C\_R from (1.3) and rho\_Upsilon' := min\{rho\_T, rho\_id, rho\_Delta, rho\_2, rho\_3, (2*C\_R)\textasciicircum{}(-1)\}, for 0 <= eta <= rho\_Upsilon' there exist an integer m >= 1, nonzero finite-dimensional Hilbert spaces L\_j and E\_j for 1 <= j <= m, a finite-dimensional Hilbert space F, operators W\_j:H->L\_j tensor E\_j, an isometry V:H->H tensor F, nonempty finite index sets Sigma\_j, unitaries U\_js on L\_j and weights p\_js for s in Sigma\_j, positive contractions C\_j on E\_j, unit vectors xi\_j in E\_j, contractions Lambda\_j:L\_j->H tensor F, CP maps Upsilon'\_j:B(H)->B(L\_j), and a CP map Upsilon':B(H)->B such that B=direct-sum\_\{j=1\}\textasciicircum{}m B(L\_j), sum\_j W\_j\textasciicircum{}*W\_j=I\_H, Delta(X)=sum\_j W\_j\textasciicircum{}*(X\_j tensor I\_Ej)W\_j for every X=(X\_1,...,X\_m) in B, Phi(T)=V\textasciicircum{}*(T tensor I\_F)V for every T in B(H), p\_js >= 0, sum\_\{s in Sigma\_j\}p\_js=1, sum\_\{s in Sigma\_j\}p\_js*(U\_js\textasciicircum{}* tensor I\_Ej)Z(U\_js tensor I\_Ej)=I\_Lj tensor (Tr\_Lj(Z)/dim(L\_j)) for every Z in B(L\_j tensor E\_j), where Tr\_Lj denotes the unnormalized partial trace over L\_j, R\_j:=sum\_\{s in Sigma\_j\}p\_js*(U\_js\textasciicircum{}* tensor I\_Ej)W\_jW\_j\textasciicircum{}*(U\_js tensor I\_Ej)=I\_Lj tensor C\_j, 0 <= C\_j <= I\_Ej, ||C\_j|| >= 1-C\_R*eta, 1-<xi\_j,C\_j\textasciicircum{}*C\_j xi\_j> <= 2*C\_R*eta, if hat-U\_js denotes U\_js in the j-th block of B and zero in every other block and iota\_js(z):=U\_js z tensor xi\_j then Lambda\_j:=sum\_\{s in Sigma\_j\}p\_js*(Delta(hat-U\_js\textasciicircum{}*) tensor I\_F)V W\_j\textasciicircum{}*iota\_js, Upsilon'\_j(T):=Lambda\_j\textasciicircum{}*(Phi(T) tensor I\_F)Lambda\_j, Upsilon'(T):=(Upsilon'\_j(T))\_\{j=1\}\textasciicircum{}m, and ||Upsilon'||\_cb <= 1. \steptag{validated} \\[6pt]
\step{1.1} Under the root hypotheses there are m >= 1 and nonzero finite-dimensional Hilbert spaces L\_j such that B=direct-sum\_\{j=1\}\textasciicircum{}m B(L\_j); complete positivity of the UCP map Delta gives finite-dimensional Hilbert spaces E\_j, provisionally allowed to be zero, and operators W\_j:H->L\_j tensor E\_j satisfying sum\_j W\_j\textasciicircum{}*W\_j=I\_H and Delta(X)=sum\_j W\_j\textasciicircum{}*(X\_j tensor I\_Ej)W\_j for every X=(X\_1,...,X\_m) in B; complete positivity and unitality of Phi give a finite-dimensional Hilbert space F and an isometry V:H->H tensor F satisfying Phi(T)=V\textasciicircum{}*(T tensor I\_F)V for every T in B(H). \steptag{validated} \\[6pt]
\step{1.1.1} Because eta <= rho\_Upsilon' <= rho\_Delta, lem-routef-delta-normalization-closeness applies to the fixed Delta' and supplies the displayed Delta as a UCP map B->B(H); lem-routef-raw-factor-setting-formation supplies B as a finite-dimensional unital C*-algebra and supplies Phi:B(H)->B(H) as the fixed UCP input map. \steptag{validated} \\[4pt]
\step{1.1.2} By GT-kitaev-fd-cstar-structure, the finite-dimensional unital C*-algebra B is *-isomorphic to direct-sum\_\{j=1\}\textasciicircum{}m B(L\_j) for some integer m>=1 and nonzero finite-dimensional Hilbert spaces L\_j; transport Delta through this *-isomorphism and henceforth identify B with that direct sum. \steptag{validated} \\[4pt]
\step{1.1.3} Apply GT-kitaev-canonical-stinespring to Delta to obtain a finite-dimensional Stinespring space G, an isometry W:H->G, and a unital *-representation pi:B->B(G) with Delta(X)=W\textasciicircum{}*pi(X)W. For each block central unit e\_j put G\_j=pi(e\_j)G, choose matrix units e\_rs\textasciicircum{}(j) in B(L\_j), put E\_j=pi(e\_11\textasciicircum{}(j))G\_j (possibly zero), and define Q\_j:L\_j tensor E\_j->G\_j by Q\_j(e\_r tensor z)=pi(e\_r1\textasciicircum{}(j))z. The matrix-unit identities show that Q\_j is unitary and Q\_j\textasciicircum{}*pi(X)Q\_j=X\_j tensor I\_Ej on G\_j. Thus W\_j:=Q\_j\textasciicircum{}*pi(e\_j)W satisfies sum\_j W\_j\textasciicircum{}*W\_j=I\_H and Delta(X)=sum\_j W\_j\textasciicircum{}*(X\_j tensor I\_Ej)W\_j. \steptag{validated} \\[4pt]
\step{1.1.4} Apply GT-kitaev-canonical-stinespring to Phi and use the same matrix-unit construction for the single full matrix algebra B(H): its finite-dimensional unital representation space is unitarily H tensor F for a finite-dimensional Hilbert space F, the representation is T->T tensor I\_F, and the Stinespring isometry becomes V:H->H tensor F, giving Phi(T)=V\textasciicircum{}*(T tensor I\_F)V. \steptag{validated} \\[4pt]
\step{1.2} For every j, let d\_j:=dim(L\_j), fix an orthonormal basis (e\_j[k])\_\{0 <= k < d\_j\}, put Sigma\_j:=\{0,...,d\_j-1\}\textasciicircum{}2, and for s=(a,b) in Sigma\_j define U\_js e\_j[k]:=exp(2*pi*i*a*k/d\_j)e\_j[(k+b) mod d\_j] and p\_js:=d\_j\textasciicircum{}(-2); then every U\_js is unitary, p\_js >= 0, sum\_\{s in Sigma\_j\}p\_js=1, and for every Z in B(L\_j tensor E\_j), sum\_\{s in Sigma\_j\}p\_js*(U\_js\textasciicircum{}* tensor I\_Ej)Z(U\_js tensor I\_Ej)=I\_Lj tensor (Tr\_Lj(Z)/d\_j), where Tr\_Lj denotes the unnormalized partial trace over L\_j; this twirl is a positive order-preserving contraction onto I\_Lj tensor B(E\_j). \steptag{validated} \\[6pt]
\step{1.2.1} For d\_j>=1, Sigma\_j=\{0,...,d\_j-1\}\textasciicircum{}2 is finite and nonempty, p\_js=d\_j\textasciicircum{}(-2) is nonnegative, and its d\_j\textasciicircum{}2 values sum to one. The formula U\_j(a,b)e\_j[k]=omega\_j\textasciicircum{}(a*k)e\_j[k+b mod d\_j], omega\_j=exp(2*pi*i/d\_j), sends the chosen orthonormal basis bijectively to an orthonormal basis, so every U\_j(a,b) is unitary. \steptag{validated} \\[4pt]
\step{1.2.2} Writing Z=sum\_\{r,t=0\}\textasciicircum{}\{d\_j-1\} E\_rt tensor Z\_rt, direct calculation gives U\_j(a,b)\textasciicircum{}* E\_rt U\_j(a,b)=omega\_j\textasciicircum{}(a*(t-r)) E\_\{r-b,t-b\}; summing first over a annihilates r!=t, and summing over b sends each diagonal E\_rr to I\_Lj. Hence the probability average is I\_Lj tensor (sum\_r Z\_rr/d\_j)=I\_Lj tensor (Tr\_Lj(Z)/d\_j). Being an average of unitary conjugations, the twirl is positive, order preserving, unital and contractive; the formula fixes exactly I\_Lj tensor B(E\_j), so it is a projection onto that subalgebra. \steptag{validated} \\[4pt]
\step{1.3} For every j define R\_j:=sum\_\{s in Sigma\_j\}p\_js*(U\_js\textasciicircum{}* tensor I\_Ej)W\_jW\_j\textasciicircum{}*(U\_js tensor I\_Ej); then R\_j=I\_Lj tensor C\_j for a positive contraction C\_j on E\_j, and, writing hat-U\_js for U\_js in the j-th block of B and zero elsewhere, W\_j\textasciicircum{}*R\_jW\_j=sum\_\{s in Sigma\_j\}p\_js*Delta(hat-U\_js\textasciicircum{}*)Delta(hat-U\_js), so lem-routef-degree-two-estimate, lem-routef-raw-factor-norms, and lem-routef-delta-normalization-closeness give ||C\_j|| >= 1-(C\_V+C\_Delta+C\_2)*eta=1-C\_R*eta. \steptag{validated} \\[6pt]
\step{1.3.1} From sum\_k W\_k\textasciicircum{}*W\_k=I\_H, each W\_j\textasciicircum{}*W\_j<=I\_H and hence ||W\_j||<=1 and 0<=W\_jW\_j\textasciicircum{}*<=I. Applying the positive unital twirl of node 1.2 to W\_jW\_j\textasciicircum{}* gives 0<=R\_j<=I and the exact formula R\_j=I\_Lj tensor C\_j with C\_j:=Tr\_Lj(W\_jW\_j\textasciicircum{}*)/d\_j, so 0<=C\_j<=I\_Ej and ||R\_j||=||C\_j||. \steptag{validated} \\[4pt]
\step{1.3.2} For hat-U\_js supported in block j, the Choi formula of node 1.1 gives Delta(hat-U\_js)=W\_j\textasciicircum{}*(U\_js tensor I\_Ej)W\_j and Delta(hat-U\_js\textasciicircum{}*)=W\_j\textasciicircum{}*(U\_js\textasciicircum{}* tensor I\_Ej)W\_j. Multiplying, averaging, and using the definition of R\_j yields the exact identity W\_j\textasciicircum{}*R\_jW\_j=sum\_s p\_js*Delta(hat-U\_js\textasciicircum{}*)Delta(hat-U\_js). \steptag{validated} \\[4pt]
\step{1.3.3} Let e\_j be the unit of the j-th block, so ||e\_j||=1. Since eta<=rho\_Upsilon'<=rho\_T, lem-routef-raw-factor-norms gives ||tilde-Delta(e\_j)||>=1-C\_V*eta; since eta<=rho\_Upsilon'<=rho\_Delta, lem-routef-delta-normalization-closeness gives ||Delta-tilde-Delta||\_cb<=C\_Delta*eta. Therefore ||Delta(e\_j)||>=1-(C\_V+C\_Delta)*eta=1-C\_N*eta. \steptag{validated} \\[4pt]
\step{1.3.4} For each s, hat-U\_js\textasciicircum{}*hat-U\_js=e\_j and both factors have norm one. Since eta<=rho\_Upsilon'<=rho\_2, lem-routef-degree-two-estimate and the probability average give ||Phi(W\_j\textasciicircum{}*R\_jW\_j)-Delta(e\_j)||<=C\_2*eta. The UCP map Phi is contractive, and ||W\_j||<=1, so ||Phi(W\_j\textasciicircum{}*R\_jW\_j)||<=||W\_j\textasciicircum{}*R\_jW\_j||<=||R\_j||=||C\_j||. Combining with node 1.3.3 gives ||C\_j||>=1-(C\_N+C\_2)*eta=1-(C\_V+C\_Delta+C\_2)*eta=1-C\_R*eta, using (1.3). \steptag{validated} \\[4pt]
\step{1.4} Since eta <= (2*C\_R)\textasciicircum{}(-1), node 1.3 gives ||C\_j|| >= 1/2, so E\_j is nonzero; finite-dimensional norm attainment gives a unit vector xi\_j in E\_j with ||C\_j xi\_j||=||C\_j||, and positivity of C\_j gives 0 <= 1-<xi\_j,C\_j\textasciicircum{}*C\_j xi\_j>=1-||C\_j||\textasciicircum{}2 <= 2*C\_R*eta. \steptag{validated} \\[6pt]
\step{1.4.1} The constants in (1.1)-(1.3) give C\_R>0, and 0<=eta<=(2*C\_R)\textasciicircum{}(-1) gives 0<=C\_R*eta<=1/2. Node 1.3 therefore yields ||C\_j||>=1-C\_R*eta>=1/2. An operator on the zero Hilbert space has norm zero, so E\_j is nonzero; because E\_j is finite-dimensional, the positive contraction C\_j attains its norm at some unit vector xi\_j, with ||C\_j xi\_j||=||C\_j||. \steptag{validated} \\[4pt]
\step{1.4.2} For that xi\_j, <xi\_j,C\_j\textasciicircum{}*C\_j xi\_j>=||C\_j xi\_j||\textasciicircum{}2=||C\_j||\textasciicircum{}2. Since C\_j is a contraction, this scalar is at most one, while ||C\_j||>=1-C\_R*eta and 0<=C\_R*eta<=1/2 imply 0<=1-<xi\_j,C\_j\textasciicircum{}*C\_j xi\_j>=1-||C\_j||\textasciicircum{}2<=1-(1-C\_R*eta)\textasciicircum{}2<=2*C\_R*eta. \steptag{validated} \\[4pt]
\step{1.5} For every j and s in Sigma\_j define iota\_js:L\_j->L\_j tensor E\_j by iota\_js(z):=U\_js z tensor xi\_j and define Lambda\_j:=sum\_\{s in Sigma\_j\}p\_js*(Delta(hat-U\_js\textasciicircum{}*) tensor I\_F)V W\_j\textasciicircum{}*iota\_js; each iota\_js and V is an isometry, Delta is UCP, ||W\_j|| <= 1, and the p\_js form a probability distribution, hence ||Lambda\_j|| <= 1. \steptag{validated} \\[6pt]
\step{1.5.1} For each j,s, the map iota\_js:z->U\_js z tensor xi\_j is an isometry because U\_js is unitary and xi\_j is a unit vector. The composition (Delta(hat-U\_js\textasciicircum{}*) tensor I\_F)V W\_j\textasciicircum{}*iota\_js is therefore correctly typed from L\_j to H tensor F, so its finite probability average defines Lambda\_j with the formula in node 1.5. \steptag{validated} \\[4pt]
\step{1.5.2} The UCP map Delta is contractive, ||hat-U\_js\textasciicircum{}*||=1, V and iota\_js are isometries, and ||W\_j\textasciicircum{}*||=||W\_j||<=1; hence every summand defining Lambda\_j has norm at most one. Since p\_js>=0 and sum\_s p\_js=1, the triangle inequality gives ||Lambda\_j||<=sum\_s p\_js=1, so Lambda\_j is a contraction. \steptag{validated} \\[4pt]
\step{1.6} Define Upsilon'\_j(T):=Lambda\_j\textasciicircum{}*(Phi(T) tensor I\_F)Lambda\_j and Upsilon'(T):=(Upsilon'\_j(T))\_\{j=1\}\textasciicircum{}m; tensoring the CP map Phi with I\_F, compression by Lambda\_j, and finite direct sums preserve complete positivity, so every Upsilon'\_j and Upsilon' are CP, and at every amplification the direct-sum maximum norm and ||Lambda\_j|| <= 1 give ||Upsilon'||\_cb <= 1 with no block-count or amplification factor. \steptag{validated} \\[6pt]
\step{1.6.1} For each j, T->Phi(T) tensor I\_F is CP because Phi is CP and tensoring with an identity representation preserves complete positivity; compression A->Lambda\_j\textasciicircum{}*A Lambda\_j is CP. Their composition Upsilon'\_j is CP. At every matrix level, positivity in the finite direct sum B=direct-sum\_j B(L\_j) is coordinatewise, so Upsilon'=(Upsilon'\_j)\_j is CP. \steptag{validated} \\[4pt]
\step{1.6.2} For every q>=1 and T in M\_q(B(H)), the amplification Phi\_q is UCP and hence contractive, and compression gives ||(Upsilon'\_j)\_q(T)||<=||Lambda\_j||\textasciicircum{}2*||Phi\_q(T)||<=||T||. The norm in M\_q(B)=direct-sum\_j M\_q(B(L\_j)) is max\_j of the component norms; therefore ||Upsilon'\_q(T)||<=||T|| for every q, and taking suprema gives ||Upsilon'||\_cb<=1 without any factor depending on q or m. \steptag{validated} \\[4pt]
\end{proofbox}

\end{document}
